\section{Objective and Research Contribution}
The approved capstone brief proposes an identity and knowledge ingestion engine,
pedagogical alignment, and student access with instructor oversight.\footnote{NUS
MComp project allocation brief, \emph{Architecting an Digital Twin for Scalable,
Style-Aligned Instructor Presence}, supplied by the student; academic advisor:
Lek Hsiang Hui. Unpublished project description.}
The intended experience is that an instructor configures course materials,
teaching preferences and learning objectives, then the Twin answers students
and initiates permitted support between their questions. Students' responses
inform subsequent support, while the instructor can review and revise the course
configuration.

The central engineering question is: \emph{How can instructor configuration,
course evidence and accumulated learner observations be connected to an agent
that initiates and continues support, and which component combinations are
supported by comparative evidence?} Here, autonomy means starting a permitted
action without a new student request or approval of every message. It includes
deciding to wait.

\begin{table}[htbp]
\centering\small
\caption{Requirements linking the project brief to executable behaviour.}
\label{tab:requirements}
\begin{tabularx}{\linewidth}{@{}L{0.29\linewidth}Y@{}}
\toprule User need & Required behaviour \\\midrule
Configure a course Twin & Review permitted sources, teaching settings and objectives; preview and publish a course-scoped release. \\
Obtain grounded support & Answer with complete, traceable course evidence, or choose clarification, abstention or refusal as appropriate. \\
Follow teaching preferences & Use the approved assistance sequence, explanation style and permitted actions in the delivered response. \\
Receive continuing support & Preserve learner observations and goals; initiate support, incorporate the response, and stop each goal using relevant evidence. \\
Retain human control & Honour membership and consent, timing limits and publication changes; support review, pause and restart. \\
\bottomrule
\end{tabularx}
\end{table}

The report makes three contributions. First, it describes a working integration
of instructor configuration, evidence-grounded dialogue and persistent proactive
tutoring. Second, it compares evidence interfaces, action planners, model
allocations and learner/timing combinations, identifying both useful changes
and unsuccessful alternatives. Third, it uses integrated trials and a reproduced
goal-completion defect to show why successful execution, correct state transitions
and useful teaching require separate evidence. The \href{https://github.com/horiiiiii032929/digital-twin}{project
repository} retains implementation and evaluation records~\cite{horinouchi2026digitaltwin}.

\subsection{Related Work and Positioning}
Retrieval-augmented generation (RAG) combines generation with retrieved external knowledge~\cite{lewis2020retrieval}.
ALCE distinguishes answer correctness from citation quality~\cite{gao2023citations}.
The implemented system enforces course/release authority and an explicit evidence-to-answer
contract. Finding a relevant passage is insufficient when a question requires
several facts or the selected source region is wrong.

StratL represents a multi-turn teaching strategy through transitions between
tutoring intents~\cite{puech-etal-2025-towards}; ScaffoldLM combines a stepwise plan
with assessment-driven memory~\cite{li-etal-2026-planning}. This project integrates those concerns with instructor-reviewed settings,
proactive execution and inspectable state.

Mixed-initiative interaction considers uncertainty, the benefit of assistance
and the cost of interrupting a user~\cite{horvitz1999mixed}. That framing motivates
comparing when to act, wait or decline, alongside what to say. The present
planning scores are engineering heuristics; their coefficients are not measured
educational effects.

MRBench evaluates tutoring through multiple pedagogical dimensions
\cite{maurya-etal-2025-unifying}. It motivates inspecting actual delivered
guidance instead of accepting an action label as proof of good instruction.
The comparisons evaluate local candidates under shared conditions.
